\section{Worker C: pole-free symbolic summation certificates}
\label{sec:sum}

\subsection{Why a rational certificate is not the whole proof}

A Gosper or creative-telescoping calculation can produce a rational function
that makes a summand telescope away from its poles. The exceptional points may
be exactly the endpoints of the finite sum. Cancelling symbolic denominators
and then substituting those endpoints is not a proof of the original sum.
SymPy provides a useful Gosper implementation for exploration, while more
general telescoping algorithms and Ore-algebra libraries supply broader
proposal mechanisms~\cite{sympy-concrete,brochet-salvy,ore-algebra}. Forge should
accept the resulting identity only with its domain and boundary obligations.

The implemented worker deliberately chooses a smaller domain in which the
certificate can be made polynomial from the outset. It handles
\begin{equation}\label{eq:sum-family}
 S_{m,W}(n)=\sum_{k=0}^{n}W(k)\binom nk^m,
 \qquad m\ge1,\quad W\in\Q[k].
\end{equation}
The checker accepts powers $1\le m\le6$ within its resource limits. The recorded
successful searches use $m=1,2$ and polynomial moments $W(k)=k^j$.
The family does not presently include $W(n,k)$: that extension requires shifting
its first argument in the $n+1$ summand, and cannot be simulated by keeping the
same $W$ on both sides.

\subsection{A boundary-safe binomial basis}

For nonnegative integers $n,k$, interpret the following quantities in $\Q$:
\[
 B(n,k)=\binom nk,
 \qquad
 A(n,k)=
 \begin{cases}0,&k=0,\\\binom n{k-1},&k>0.
 \end{cases}
\]
As usual $\binom nk=0$ for $k>n$. Two identities hold on the entire nonnegative
index domain:
\begin{align}
 \binom{n+1}{k}&=A(n,k)+B(n,k),\label{eq:pascal-safe}\\
 (n+1-k)A(n,k)-kB(n,k)&=0.\label{eq:cross-safe}
\end{align}
Here $n+1-k$ is subtraction after embedding in $\Q$, not truncated subtraction
in $\N$. For $k>n+1$, both binomial factors vanish. At $k=0$, the separate
zero definition of $A$ is essential. Defining $A(n,0)$ using
\code{Nat.choose n (0 - 1)} would give $\binom n0=1$, because natural
subtraction truncates. That is precisely the wrong endpoint.

The pinned Mathlib source provides \code{Nat.choose}, the successor/Pascal
identity, and vanishing above the top index~\cite{mathlib-choose}. A Lean bridge
should prove \eqref{eq:pascal-safe} and \eqref{eq:cross-safe} once, splitting
$k=0$, $1\le k\le n+1$, and $k>n+1$ as needed. These primitive theorems can then
be reused for every certificate in the family.

\subsection{The exact certificate language}

Search for $p_0,p_1\in\Q[n]$, a polynomial $r\in\Q[n,k]$, and a polynomial
$h\in\Q[A,B,n,k]$ such that
\begin{align}
 &p_1(n)W(k)(A+B)^m+p_0(n)W(k)B^m \notag\\[-1mm]
 &\quad-r(n,k+1)B^m+r(n,k)A^m
 =h(A,B,n,k)\big((n+1-k)A-kB\big),\label{eq:sum-cert}
\end{align}
with $p_1\not\equiv0$. Every symbol here is a polynomial indeterminate. In
particular, the coefficient field is $\Q$, not $\Q(n,k)$. Division by a
polynomial in $n$ or $k$ is not allowed to hide inside the coefficient domain.

The certificate stores the four polynomials $p_0,p_1,r,h$. The problem supplies
$m$ and $W$. The checker independently expands both sides of
\eqref{eq:sum-cert} by sparse rational arithmetic. It does not need to run
Gosper's algorithm, test a rational-function identity, or trust a simplifier's
handling of singular values.

\begin{theorem}[Pole-free sum recurrence]\label{thm:sum}
If \eqref{eq:sum-cert} holds for $m\ge1$, then
\begin{equation}\label{eq:sum-rec}
 p_1(n)S_{m,W}(n+1)+p_0(n)S_{m,W}(n)=0
 \qquad(n\in\N).
\end{equation}
\end{theorem}
\begin{proof}
Substitute $A=A(n,k)$ and $B=B(n,k)$ into the polynomial certificate. The
right-hand side vanishes by \eqref{eq:cross-safe}. Define
\[
 T(n,k)=r(n,k)A(n,k)^m.
\]
Since $A(n,k+1)=B(n,k)$, the remaining identity is
\[
 p_1(n)W(k)\binom{n+1}k^m+
 p_0(n)W(k)\binom nk^m=T(n,k+1)-T(n,k).
\]
Sum this equality for $k=0,\ldots,n+1$. The first term is exactly
$p_1(n)S_{m,W}(n+1)$. The extra $k=n+1$ term from the second row is zero, so the
second term is $p_0(n)S_{m,W}(n)$. The right side is
$T(n,n+2)-T(n,0)$. Both endpoints vanish: $A(n,0)=0$ and
$A(n,n+2)=\binom n{n+1}=0$, and $m\ge1$. This proves the recurrence without
any denominator cancellation.
\end{proof}

The condition $m\ge1$ is mathematical, not just a convenient validation rule.
For $m=0$ the endpoint factors would be $0^0=1$ in the usual algebraic power
convention, so the argument would not establish zero boundary terms. The checker
rejects that input rather than silently reusing the theorem.

\subsection{Synthesis by a finite nullspace computation}

Choose degree bounds $d_p,d_r$ and the templates
\[
 p_i(n)=\sum_{j=0}^{d_p}c_{ij}n^j,\qquad
 r(n,k)=\sum_{a+b\le d_r}d_{ab}n^ak^b.
\]
There are
\[
 N=2(d_p+1)+\binom{d_r+2}{2}
\]
unknown coefficients. The residual on the left of \eqref{eq:sum-cert} is
linear in these unknowns.

Reduce each coefficient column modulo the principal ideal generated by
$\rho=(n+1-k)A-kB$ in $\Q[A,B,n,k]$. Collect the coefficients of all remaining
monomials into an exact rational matrix $M$. A vector in $\ker M$ describes a
residual that lies in $\Ideal\rho$. Retain a vector only if its $p_1$ component
is nonzero. Multiply by a common denominator, remove the integer content, and
choose a positive leading coefficient of $p_1$ for a deterministic display.
Finally, divide the completed residual by $\rho$ to obtain $h$ and replay
\eqref{eq:sum-cert} independently.

\begin{proposition}[Fixed-template coverage]
Exact nullspace computation finds a nontrivial order-one certificate of this
form whenever one exists within the chosen $p_i$ and $r$ templates.
\end{proposition}
\begin{proof}
Normal-form reduction by a fixed Gr\"obner basis is $\Q$-linear. Thus the matrix
kernel is exactly the coefficient vectors whose residual has zero normal form,
which is exactly membership in the principal ideal. The projection of a kernel
basis onto the $p_1$ coefficients is nonzero precisely when that projection is
nonzero for some kernel vector. Hence inspecting a full kernel basis suffices
to find a permitted nonzero $p_1$ component.
\end{proof}

The synthesis loop tries degree pairs in increasing total template cost. It is
not a general creative-telescoping decision procedure: it fixes recurrence
order one, a particular pole-free certificate form, and finite degree bounds.
A missing certificate means \Unknown for the sum identity, even if the finite
linear template search was fully exhausted. The successful larger-template
runs and the unsuccessful smaller-template runs are both retained in the
experimental corpus.

\subsection{Worked example: the squared-binomial identity}
\label{sec:binomial-square}

For $m=2$, $W=1$, the search produces
\[
 p_1=n+1,\qquad p_0=-2(2n+1),\qquad r=2k-3n-3,
 \qquad h=2(B-A).
\]
The complete local certificate is
\begin{align*}
 &(n+1)(A+B)^2-2(2n+1)B^2\\
 &\quad-\bigl(2(k+1)-3n-3\bigr)B^2
       +(2k-3n-3)A^2\\
 &=2(B-A)\bigl((n+1-k)A-kB\bigr).
\end{align*}
It can be verified by an ordinary ring normalization. Theorem~\ref{thm:sum}
then gives
\[
 (n+1)S(n+1)=2(2n+1)S(n),\qquad S(0)=1.
\]
The sequence $H(n)=\binom{2n}{n}$ has the same initial value and recurrence:
for example its factorial ratio gives
\[
 \frac{H(n+1)}{H(n)}=
 \frac{(2n+2)(2n+1)}{(n+1)^2}
 =\frac{2(2n+1)}{n+1},
\]
where the positive integer denominators are nonzero. Since $n+1\ne0$ on
$\N$, induction establishes
\[
 \boxed{\displaystyle\sum_{k=0}^n\binom nk^2=\binom{2n}{n}.}
\]
This last comparison is a mathematical derivation in the article. The executed
Python certificate certifies the recurrence, not a Lean theorem about the closed
form.

A rational Gosper representation of the same telescoping term is
\[
 \frac{k^2(2k-3n-3)}{(n+1-k)^2}\binom nk^2.
\]
Its denominator vanishes at $k=n+1$. Rewriting the \emph{entire term} as
$(2k-3n-3)A(n,k)^2$ gives a value on that boundary and permits a direct pointwise
proof there. This is not the assertion that two partial rational expressions
agree at a pole. It is a new total expression accompanied by an identity valid
at every required index.

For $m=1$, $W=1$, the search finds $p_1=1$, $p_0=-2$, $r=-1$ and the primitive
relation multiplier $h=0$. The corresponding recurrence is $S(n+1)=2S(n)$;
with $S(0)=1$ it gives the usual $2^n$ sum.

\subsection{Initial values and singular leading coefficients}
\label{sec:singular}

A recurrence is not a uniqueness theorem without a nonvanishing condition or
additional seeds. The first-moment case produces
\[
 nS(n+1)-2(n+1)S(n)=0.
\]
At $n=0$, this equation forces $S(0)=0$ but says nothing about $S(1)$. The true
sequence $S(n)=\sum_k k\binom nk$ and the wrong sequence $2S(n)$ satisfy the
same recurrence and the same $n=0$ seed. The regression suite explicitly checks
this counterexample to a base-only acceptance rule.

\begin{proposition}[Order-one recurrence comparison]\label{prop:seeds}
Suppose $u,v$ satisfy the same recurrence
$a(n)z(n+1)+b(n)z(n)=0$ on $\N$. Let
$E=\{e\in\N:a(e)=0\}$. If $u(0)=v(0)$ and
$u(e+1)=v(e+1)$ for every $e\in E$, then $u=v$ on $\N$.
\end{proposition}
\begin{proof}
Induct on $n$. To obtain equality at $n+1$, use the corresponding extra seed
when $a(n)=0$. Otherwise subtract the two recurrence equations, use equality at
$n$, and cancel the nonzero scalar $a(n)$.
\end{proof}

For nonzero rational-polynomial $a$, $E$ is finite. A production comparator can
clear constant denominators, remove the power of $n$, and enumerate possible
integer roots using divisibility of the constant coefficient. Each proposed
root is checked by evaluation, and the divisibility argument accounts for all
remaining roots. A simpler but potentially expensive checker can use a Cauchy
root bound and evaluate every integer below it. Either route requires a proof
of coverage; a CAS list of roots is not itself that proof.

The prototype records nonnegative roots of $p_1$ as \emph{search diagnostics}.
The independent recurrence checker does not use those diagnostics to accept a
closed form. The Lean recurrence-comparison layer and certified seed-coverage
check are proposed work, not implemented artifacts.

\subsection{The seven synthesized recurrences}

Table~\ref{tab:recurrences} reports the actual recurrences, with
$S_{m,j}(n)=\sum_{k=0}^n k^j\binom nk^m$ and $k^0=1$. They all have the form
$p_1(n)S_{m,j}(n+1)+p_0(n)S_{m,j}(n)=0$. Common factors are not silently
cancelled, because cancellation could alter the exceptional-index obligations.

\begingroup\small
\begin{longtable}{@{}ccp{0.27\linewidth}p{0.28\linewidth}c@{}}
\caption{Synthesized recurrence coefficients and total degree of $r(n,k)$, generated from the saved certificates.}\label{tab:recurrences}\\
\toprule
$m$ & $j$ & $p_1(n)$ & $p_0(n)$ & $\deg r$ \\
\midrule\endhead
1 & 0 & $1$ & $-2$ & 0 \\
1 & 1 & $n$ & $- 2 \left(n + 1\right)$ & 2 \\
1 & 2 & $n$ & $- 2 \left(n + 2\right)$ & 3 \\
1 & 3 & $n^{2} \left(n + 3\right)$ & $- 2 \left(n + 1\right)^{2} \left(n + 4\right)$ & 6 \\
2 & 0 & $n + 1$ & $- 2 \left(2 n + 1\right)$ & 1 \\
2 & 1 & $n^{2}$ & $- 2 n \left(2 n + 1\right)$ & 3 \\
2 & 2 & $n^{3}$ & $- 2 \left(n + 1\right)^{2} \left(2 n - 1\right)$ & 5 \\
\bottomrule
\end{longtable}

\endgroup


The two highest-moment searches fail under $d_r\le4$ and succeed when the
explicitly stated bound is enlarged. An order-one attempt for $m=3,W=1$ remains
\Unknown at the recorded small bounds. No claim is made that the sequence lacks
a higher-order recurrence or that the search has exhausted creative telescoping
in general.

\subsection{A route to general creative telescoping}

The next interface should accept a telescoper
$L=\sum_{i=0}^\ell p_i(n)E_n^i$, a finite set of pointwise summand relations,
and a certificate for
\[
 \sum_{i=0}^\ell p_i(n)f(n+i,k)=G(n,k+1)-G(n,k).
\]
The acceptance pipeline must also prove a common summation range, any added or
removed tail terms, endpoint values of $G$, and enough initial data to compare
the resulting recurrence with the target sequence. Nonzero denominators become
explicit subgoals; singular indices become separate branches. When endpoints
are not zero, their exact difference belongs on the right-hand side of the
recurrence rather than being discarded.

A Sage \code{ore\_algebra} process is a concrete optional proposal backend:
its documented functionality includes creative telescoping, desingularization,
and operations on D-finite descriptions~\cite{ore-algebra}. The compact
certificate methods of Brochet and Salvy provide an additional algorithmic
reference~\cite{brochet-salvy}. Neither is a Lean oracle. Their output still has
to be translated into the pointwise relations and boundary certificate accepted
by Forge. This interface is designed here, whereas the delivered code implements
only the polynomial binomial template.
